% Partie : sets-functions-relations
% Chapitre : functions
% Section : basics

\documentclass[../../../include/open-logic-section]{subfiles}

\begin{document}

\olfileid[fr]{sfr}{fun}{bas}
\olsection{Notions de base}

\begin{explain}
Une \emph{fonction} associe à chaque !!{element} d’un ensemble donné
un !!{element} déterminé d’un autre ensemble donné, qui peut être le même.
Par exemple, l’opération qui consiste à ajouter~$1$ définit une fonction :
à chaque nombre~$n$, elle associe l’unique nombre~$n+1$.

Plus généralement, une fonction peut recevoir en entrée des couples,
des triplets, etc., et produire une sortie. L’arithmétique élémentaire
nous fournit de nombreuses fonctions familières : l’addition et la
multiplication, par exemple, reçoivent deux nombres et en renvoient un troisième.

Dans ce sens mathématique abstrait, une fonction est une \emph{boîte noire} :
seule compte la sortie associée à chaque entrée, et non la méthode
qui permet de la calculer.
\end{explain}

\begin{defn}[Fonction]
Une \emph{fonction} $f \colon A \to B$ associe à chaque !!{element}
de~$A$ un unique !!{element} de~$B$.

On appelle $A$ le \emph{domaine} de~$f$ et $B$ son \emph{ensemble d’arrivée}.
Les !!{element}s de~$A$ sont les entrées ou \emph{arguments} de~$f$.
L’!!{element} de~$B$ que~$f$ associe à un argument~$x$ est appelé
la \emph{valeur de~$f$} en~$x$ et se note~$f(x)$.

L’\emph{image} $\ran{f}$ de~$f$ est la partie de son ensemble d’arrivée
constituée des valeurs effectivement prises par~$f$ :
$\ran{f} = \Setabs{f(x)}{x \in A}$.
\end{defn}

Le diagramme de la \olref{fig:function} aide à se représenter une fonction.
L’ellipse de gauche représente son \emph{domaine}, celle de droite
son \emph{ensemble d’arrivée}. Chaque flèche part d’un \emph{argument}
du domaine et aboutit à la \emph{valeur} correspondante dans l’ensemble d’arrivée.

\begin{figure}
  \olasset{assets/diagrams/function.tikz}
  \caption{Une fonction associe à chaque !!{element} d’un ensemble
    !!a{element} d’un autre ensemble. Une flèche part d’un argument
    du domaine et aboutit à la valeur correspondante dans l’ensemble d’arrivée.}
  \ollabel{fig:function}
\end{figure}

\begin{ex}
La multiplication reçoit en entrée des couples d’entiers naturels et
renvoie des entiers naturels. Elle va donc de $\Nat \times \Nat$
(son domaine) dans $\Nat$ (son ensemble d’arrivée).
Son image est elle aussi $\Nat$, puisque tout $n \in \Nat$ est égal à
$n \times 1$.
\end{ex}

\begin{ex}
La multiplication est une fonction parce qu’elle associe à chaque entrée,
c’est-à-dire à chaque couple d’entiers naturels, une unique sortie :
$\times \colon \Nat^2 \to \Nat$. En revanche, associer à un entier
naturel~$n$ un réel~$x$ tel que $x^2 = n$ ne définit pas une fonction :
chaque entier strictement positif~$n$ possède deux racines carrées,
$\sqrt{n}$ et~$-\sqrt{n}$. On obtient une fonction en ne retenant que
la racine carrée positive ou nulle :
$\sqrt{\phantom{X}} \colon \Nat \to \Real$.
\footnote{Précision éditoriale : on inclut la racine nulle pour que
la fonction soit aussi définie en zéro, qui appartient à son domaine.}
\end{ex}

\begin{ex}
La relation qui associe à chaque étudiant d’un cours sa note finale est
une fonction : un étudiant ne peut pas recevoir deux notes finales
différentes pour un même cours. La relation qui associe à chaque étudiant
ses parents n’est pas une fonction : un étudiant peut avoir zéro, deux
ou davantage de parents.
\end{ex}

\begin{explain}
Pour définir une fonction, on peut préciser sa valeur pour chaque
argument possible. On peut le faire par une formule, par une méthode
de calcul ou par une liste des valeurs correspondant aux arguments.
Quelle que soit la méthode choisie, il faut s’assurer qu’à chaque
argument correspond une valeur et une seule.
\end{explain}

\begin{ex}
Soit $f \colon \Nat \to \Nat$ définie par $f(x) = x+1$.
Cette définition précise que $f$ reçoit des entiers naturels et renvoie
des entiers naturels : à tout entier naturel~$x$, elle associe son
successeur~$x+1$. Ici, l’ensemble d’arrivée $\Nat$ n’est pas l’image
de~$f$, car l’entier naturel~$0$ n’est le successeur d’aucun entier naturel.
L’image de~$f$ est l’ensemble des entiers strictement positifs, $\Int^{+}$.
\end{ex}

\begin{ex}\ollabel{examplefunext}
Soit $g \colon \Nat \to \Nat$ définie par $g(x) = x+2-1$.
C’est une fonction qui reçoit des entiers naturels et renvoie des
entiers naturels. À un entier naturel~$x$, elle associe le prédécesseur
du successeur du successeur de~$x$, soit $x+1$.
\footnote{Correction éditoriale : l’original appelait l’entrée
$n$ avant de la désigner par $x$ ; on emploie ici la même lettre.}
\end{ex}

\begin{explain}
Nous venons de considérer deux fonctions $f$ et $g$ données par des
\emph{définitions} différentes. Pourtant, il s’agit de la \emph{même fonction}.
Pour tout entier naturel~$n$, on a en effet
$f(n) = n+1 = n+2-1 = g(n)$. Autrement dit, nos définitions de $f$ et~$g$
décrivent la même correspondance au moyen d’équations différentes.
Nous utilisons donc implicitement un principe d’extensionnalité
pour les fonctions :
\[
  \text{si }\forall x\, f(x) = g(x)\text{, alors }f = g
\]
à condition que $f$ et~$g$ aient le même domaine et le même ensemble d’arrivée.
\end{explain}

\begin{ex}
On peut aussi définir une fonction par cas. Définissons par exemple
$h \colon \Nat \to \Nat$ par
\[
h(x) =
\begin{cases}
  \frac{x}{2} & \text{si $x$ est pair} \\
  \frac{x+1}{2} & \text{si $x$ est impair.}
\end{cases}
\]
Tout entier naturel étant pair ou impair, cette fonction renvoie toujours
un entier naturel. Dans une telle définition par cas disjoints, toute
entrée possible doit relever d’un cas et d’un seul. Il faut parfois
démontrer que les cas sont exhaustifs et mutuellement exclusifs.
\footnote{Précision éditoriale : l’original présente cette condition comme
une exigence générale. Des cas peuvent aussi se recouper, à condition
qu’ils attribuent la même valeur aux entrées communes ; toutes les entrées
doivent néanmoins être couvertes.}
\end{ex}

\end{document}
